\documentclass[a4paper]{article}

\usepackage[english]{babel}
\usepackage{hyperref}
\usepackage{float}
\usepackage[utf8]{inputenc}
\usepackage{amsmath}
\usepackage{graphicx}
\usepackage[colorinlistoftodos]{todonotes}
\usepackage{listings}
\usepackage{xcolor}

% Define colors for terminal-style code
\definecolor{terminalback}{RGB}{40,44,52}
\definecolor{terminaltext}{RGB}{171,178,191}
\definecolor{terminalcomment}{RGB}{92,99,112}
\definecolor{terminalkeyword}{RGB}{198,120,221}

% Terminal-style listing
\lstdefinestyle{terminal}{
    backgroundcolor=\color{terminalback},
    basicstyle=\ttfamily\small\color{terminaltext},
    breaklines=true,
    frame=single,
    frameround=tttt,
    rulecolor=\color{terminalback},
    numbers=none,
    xleftmargin=10pt,
    xrightmargin=10pt,
    aboveskip=10pt,
    belowskip=10pt,
}

% Code listing style
\lstdefinestyle{code}{
    backgroundcolor=\color{white},
    basicstyle=\ttfamily\small,
    breaklines=true,
    frame=single,
    numbers=left,
    numberstyle=\tiny\color{gray},
    keywordstyle=\color{blue},
    commentstyle=\color{green!60!black},
    stringstyle=\color{red},
    xleftmargin=15pt,
    xrightmargin=10pt,
    aboveskip=10pt,
    belowskip=10pt,
}

\title{Adding Custom System Calls to xv6-riscv Operating System}

\author{\normalsize Alireza Doostimehr \\\normalsize Student Number: 400442135}
\date{\color{black}February 2026}

\begin{document}
\maketitle

\section{Introduction}

This project adds a custom \texttt{getticks} system call to the xv6-riscv operating system. This system call returns the number of timer ticks since the system started running.

The complete source code and implementation details are available at:
\url{https://github.com/alirezadoostimehr/xv6-no-risks}

\section{Implementation Process}

Adding a system call to xv6-riscv requires changes to several files across both user space and kernel space. The following sections explain each step in detail.

\subsection{System Call Number Assignment}

Every system call needs a unique identification number. This number is used by the kernel to determine which function to execute when a system call is made.

\textbf{File:} \texttt{kernel/syscall.h}

\begin{lstlisting}[style=code, language=C]
#define SYS_getticks 23
\end{lstlisting}

This definition assigns the number 23 to our new system call. The kernel uses this number to look up the correct function in its system call table.

\subsection{User-Space Interface}

User programs need to know the function signature to call the system call properly.

\textbf{File:} \texttt{user/user.h}

\begin{lstlisting}[style=code, language=C]
int getticks(void);
\end{lstlisting}

This declaration tells the compiler that \texttt{getticks} is a function that takes no arguments and returns an integer value.

\subsection{User-Space Stub Creation}

The user-space stub generates the assembly code that performs the actual system call.

\textbf{File:} \texttt{user/usys.pl}

\begin{lstlisting}[style=code]
entry("getticks");
\end{lstlisting}

This Perl script generates assembly code that loads the system call number into register \texttt{a7} and executes the \texttt{ecall} instruction to trap into kernel mode.

\subsection{Kernel Implementation}

The kernel function implements the actual logic of the system call.

\textbf{File:} \texttt{kernel/sysproc.c}

\begin{lstlisting}[style=code, language=C]
uint64 sys_getticks(void) {
  uint temp;
  acquire(&tickslock);
  temp = ticks;
  release(&tickslock);
  return temp;
}
\end{lstlisting}

This function reads the global \texttt{ticks} variable, which is updated by the timer interrupt handler. The function uses a spinlock (\texttt{tickslock}) to ensure thread-safe access, preventing race conditions when multiple processes might try to read the ticks value at the same time.

\subsection{System Call Registration}

The kernel needs to know which function to call when it receives a system call request.

\textbf{File:} \texttt{kernel/syscall.c}

First, we declare the external function:

\begin{lstlisting}[style=code, language=C]
extern uint64 sys_getticks(void);
\end{lstlisting}

Then, we add it to the system call table:

\begin{lstlisting}[style=code, language=C]
static uint64 (*syscalls[])(void) = {
  // ... other system calls ...
  [SYS_getticks] sys_getticks,
};
\end{lstlisting}

This array maps system call numbers to their corresponding kernel functions. When a user program executes \texttt{ecall} with system call number 23, the kernel looks up this array and calls \texttt{sys\_getticks()}.

\subsection{Test Program}

To verify the system call works correctly, we created a test program.

\textbf{File:} \texttt{user/getticks.c}

\begin{lstlisting}[style=code, language=C]
#include "kernel/types.h"
#include "kernel/stat.h"
#include "user/user.h"

int main(int argc, char *argv[]) {
  printf("Now: %d\n", getticks());
  exit(0);
}
\end{lstlisting}

\subsection{Build Configuration}

Finally, we need to tell the build system to compile our code.

\textbf{File:} \texttt{Makefile}

\begin{lstlisting}[style=code]
UPROGS=\
  # ... other programs ...
  $U/_getticks\
\end{lstlisting}

\section{Execution Flow}

When a user types \texttt{getticks} in the xv6 shell and presses enter, the following sequence of events occurs:

\begin{enumerate}
    \item The shell executes the \texttt{getticks} program from the filesystem
    \item The program starts running in user mode and calls the \texttt{getticks()} function
    \item The user-space stub (generated by \texttt{usys.pl}) loads syscall number 23 into register \texttt{a7}
    \item The stub executes the \texttt{ecall} instruction, which triggers a trap into kernel mode
    \item The CPU switches to supervisor mode and jumps to the trap handler
    \item The kernel reads the syscall number from register \texttt{a7}
    \item The kernel looks up \texttt{syscalls[23]} in the system call table and finds \texttt{sys\_getticks}
    \item The kernel calls \texttt{sys\_getticks()}, which:
    \begin{itemize}
        \item Acquires the \texttt{tickslock} spinlock
        \item Reads the current value of the \texttt{ticks} global variable
        \item Releases the lock
        \item Returns the value
    \end{itemize}
    \item The kernel puts the return value in register \texttt{a0}
    \item The CPU returns to user mode
    \item The \texttt{getticks()} function returns with the value from \texttt{a0}
    \item The program prints the result using \texttt{printf}
    \item The program exits
\end{enumerate}

The \texttt{ticks} variable is a global counter in the kernel that gets incremented every time a timer interrupt occurs.

\section{Testing Results}

The \texttt{getticks} system call was tested by running the test program at different times after system boot:

\begin{lstlisting}[style=terminal]
$ getticks
Now: 271
$ getticks
Now: 329
\end{lstlisting}

The increasing tick values show that the system call correctly reads the kernel's tick counter.

\subsection{Additional Work}

A \texttt{revrev} system call was also added as a fun exercise. This system call takes a string as input, reverses it, and prints the reversed string.

\end{document}